\section{Linking differences and degree-three Tate cell attachments}
\label{sec:136-degree-three-linking-repair}

\paragraph{Archive and exact scope.}
The uniform category is cote136, physical pp.~117 and119,
footers116 and118, scans117 and119; the common-source question
is physical p.~121, footer120, scan121.  This result treats
an actual extension under an old algebra.  It preserves the
vertex-residue compatibility needed to repair degree two,
for degree-three Tate cells attached along their full simplicial
boundaries.  A separate final criterion explains why replacing
those boundaries by only the vertices introduces a further gate.
No completion of the original integral common-source problem
is asserted.

Work over \(R=\mathbb Z\) or \(\mathbb Z_{(p)}\), with
\(S=\Gamma_R(T)\), and put
\[
                         J=\mathbb Q[T]/S.
                                                               \tag{136.LK1}
\]
ARCH-136-GRADED-INJECTIVES proves that \(J\) is injective
in the category of nonnegative weight-graded \(S\)-modules.
This is the only injectivity assertion used below.  In particular
no injectivity of a simplicial sphere object is presumed.

\subsection{The hypotheses and the actual cell algebra}

Let \(A\) be a simplicial nonnegative strict differential
PD \(S\)-algebra with \(A_0=S\), and suppose, literally
in every weight and simplex,
\[
 H^0(A_n)=S,\qquad H^1(A_n)=0,
                 H^2(A_n)\text{ is integer torsion}.
                                                               \tag{136.LK2}
\]
The all-weight assertion in the first two conditions is essential.
A low-weight exception multiplied by a new generator of large
weight would otherwise enter the output range considered below.
These conditions hold in the guarded constructions when every
noncoefficient generator is born above the fixed RH1 cutoff,
the low repair is performed, and rational vertical cohomology
is the coefficient algebra.

Fix finitely many admissible targets \(C_j\), with
\((C_j)_0=S\), and strict maps \(f_j:A\to C_j\).
Assume their degree-two vertex residues agree in all weights
\(w\geq W_0\).  Explicitly, if
\[
 x\in Z^2(A_n^w),\qquad db=Nx,\qquad N\geq1,
                                                               \tag{136.LK3}
\]
choose \(a_j\in C_j^{w,1}(\Delta[n])\) with
\(da_j=f_jx\), and then \(q_j\in C_j^{w,0}(\Delta[n])\)
with \(dq_j=f_jb-Na_j\).  The vertex values of \(q_j\)
define an element of
\[
              S^{w\,\oplus(n+1)}/
                    (\operatorname{diag}S^w+N S^{w\,\oplus(n+1)}).
                                                               \tag{136.LK4}
\]
The hypothesis is equality of (136.LK4) for every target.
Enlarge \(W_0\) to include the target vertical comparison
through degree two; this makes all indicated choices available
and their residue classes well-defined, as in the checked
degree-two-residue theorem.

For each cell \(\alpha\), choose a simplicial dimension
\(n_\alpha\geq1\), a weight \(v_\alpha>0\), an integer
\(r_\alpha\geq2\), and old sections
\[
 z_\alpha\in A^{v_\alpha,3}(\Delta[n_\alpha]),\qquad
 e_\alpha\in A^{v_\alpha,2}(\Delta[n_\alpha]),\qquad
 dz_\alpha=0,\quad de_\alpha=r_\alpha z_\alpha.
                                                               \tag{136.LK5}
\]
On the entire boundary \(\partial\Delta[n_\alpha]\),
assume old sections \(c_\alpha^\partial,t_\alpha^\partial\)
of forms two and one have already been specified, satisfying
\(dc_\alpha^\partial=z_\alpha|_\partial\) and
\(dt_\alpha^\partial=e_\alpha|_\partial-r_\alpha
c_\alpha^\partial\).
Form the actual free relative PD pushout \(A\to B\)
adjoining one cell pair
\[
       |c_\alpha|=2,\quad |t_\alpha|=1,\qquad
       dc_\alpha=z_\alpha,\quad
       dt_\alpha=e_\alpha-r_\alpha c_\alpha,
                                                               \tag{136.LK6}
\]
and identifying its full boundary with that old data.  Any set
of cells may be attached in this one stage, with all boundaries
in \(A\).  At each simplex the underlying algebra is free
over \(A_n\) on the pairs indexed by relative basis simplices;
it has no other new relation.  This follows by eliminating the
boundary basis generators in the degreewise split free-module
inclusion.  In particular \(B_0=S\).

Assume strict extensions \(g_j:B\to C_j\) of the given
maps have been provided.  This hypothesis includes the actual
installation of the prescribed boundary data; the theorem below
does not supply it silently.  Finally choose one \(V\),
depending only on the targets, such that in weights \(v\geq V\)
the endpoints of a target zero-form on \(\Delta[1]\) can be
prescribed arbitrarily in \((S^v)^2\), and the target
\(H^0,H^1\) have their coefficient values.  Target row
acyclicity in form zero and vertical admissibility supply such
a \(V\).  Require \(v_\alpha\geq V\) for the one-dimensional
cells; no maximum of the cell weights is assumed.

\subsection{The exact low cohomology calculation}

Let
\[
 L=\bigoplus_\alpha(S/r_\alpha S)\langle v_\alpha\rangle
       \otimes_R
       R[\Delta[n_\alpha]]/R[\partial\Delta[n_\alpha]].
                                                               \tag{136.LK7}
\]
There is a natural simplicial map
\(\psi:L\to H^3(A)\) taking the coefficient of the new
\(c_\alpha\) to \([z_\alpha]\).  It is well-defined
because \(r_\alpha[z_\alpha]=0\), and its boundary values
are zero because the prescribed \(c_\alpha^\partial\)
are primitives there.  Write \(K=\ker\psi\).  Then
\[
 H^0(B)=S,\quad H^1(B)=0,\qquad
 0\longrightarrow H^2(A)\longrightarrow H^2(B)
                   \longrightarrow K\longrightarrow0
                                                               \tag{136.LK8}
\]
are natural assertions of simplicial weight-graded modules.

Here is a direct proof, including the mixed terms.  In one
simplex write the new variables as \(c_i,t_i\).  A closed
one-form \(a+\sum q_i t_i\) has \(dq_i=0\) and
\(r_iq_i=0\), by its independent \(t_i\) and \(c_i\)
coefficients.  Thus \(q_i\in S\), and torsion-freeness
of \(S\) gives \(q_i=0\).  The first two assertions follow.
A two-form also has possible terms \(h_{ij}t_it_j\),
\(i<j\).  In a cycle their coefficients satisfy
\(dh_{ij}=0\), from the independent terms with two new odd
variables and an old one-form.  Hence \(h_{ij}\in S\).
The coefficients of \(c_it_j\) and \(c_jt_i\) in its
differential force \(r_i h_{ij}=r_j h_{ij}=0\), so these
mixed terms are zero.  No other differential in form two
produces a new \(c_i t_j\) term.

The remaining cycle is
\(a+\sum b_i t_i+\sum q_i c_i\), and the cycle equations are
\[
 db_i=0,\qquad dq_i+r_i b_i=0,\qquad
              da-\sum b_i e_i+\sum q_i z_i=0.
                                                               \tag{136.LK9}
\]
Write \(b_i=d\beta_i\), using \(H^1(A)=0\), and
subtract \(d(\sum\beta_i t_i)\).  The resulting representative
is \(a'+\sum s_i c_i\), with \(s_i=q_i+r_i\beta_i\in S\)
and \(da'+\sum s_i z_i=0\).  Its coefficient vector modulo
the \(r_i\) lies in \(K\).  Every vector in \(K\) has
such a lift.  Changing a representative by a boundary changes
\(s_i\) by \(-r_i\lambda_i\) with \(\lambda_i\in S\).
An old two-cycle can be a new boundary only if all these
coefficients are zero; torsion-freeness then forces
\(\lambda_i=0\).  This proves (136.LK8).  The description
is equivalently the relative-cohomology projection of the
actual inclusion \(A\to B\), so it is natural under all
simplicial operators, including boundary faces that land in
the old algebra.  All sums in this calculation are finite.

In particular \(H^2(B)\) is integer torsion, with no
uniform bound on its annihilators required.

\subsection{The universal linking difference}

Put \(\mathcal V_n(M)=M^{\oplus(n+1)}/\operatorname{diag}M\),
with the simplicial maps induced by deleting or repeating
vertices.  For a torsion class \([x]\in H^2(B_n^w)\),
\(w\geq W_0\), choose (136.LK3) in \(B\) and target
primitives as above.  Define
\[
 \Delta_j([x])=
 \left[\frac{q_j|_{\mathrm{vertices}}
                     -q_1|_{\mathrm{vertices}}}{N}\right]
                   \in\mathcal V_n(J)^w.
                                                               \tag{136.LK10}
\]
Here \(\mathcal V_n(J)\) is canonically
\(\mathcal V_n(S\otimes\mathbb Q)/\mathcal V_n(S)\).
The class is independent of every choice.  Changing a target
primitive \(a_j\) changes it by an exact one-form, so the
resulting change after division by \(N\) is an integral
vertex vector.  Changing \(q_j\) by a closed zero-form
changes only a diagonal coefficient.  Changing \(b\) by a
closed source one-form changes it by \(dh\), since
\(H^1(B)=0\); the vertex values of every \(g_jh\)
agree because \(B_0=S\).  Changing the representative
\(x\) by a boundary is handled by changing \(a_j\)
by the image of its source primitive.  Passing to a common
multiple of annihilators proves independence of \(N\)
and additivity.  The same formulas prove weight-preserving
\(S\)-linearity and simplicial naturality.

Thus \(\Delta_j\) is a genuine map of simplicial graded
modules on the weight tail of \(H^2(B)\).  Its vanishing
is exactly equality of the finite-modulus residues (136.LK4)
for every \(N\): a vector divided by \(N\) is integral
modulo the diagonal precisely when its original vector is
zero modulo \(N\) and that diagonal.  By the old compatibility
hypothesis, (136.LK10) vanishes on \(H^2(A)\) in that
tail.  It consequently descends, by (136.LK8), to
\[
                         K_{\geq W_0}\longrightarrow\mathcal V(J).
                                                               \tag{136.LK11}
\]

\paragraph{ARCH-136-DEGREE-THREE-LINKING-REPAIR.}
Under the stated hypotheses, one may change only the images
of the new one-forms \(t_\alpha\) in one-dimensional cells,
by adding closed one-forms relative to their full boundaries,
so that every degree-two vertex residue of \(B\) agrees
across all the targets in the same uniform weight tail.
The maps on \(A\), all prescribed boundary data, and all
images of \(c_\alpha\) are retained.  The result covers
arbitrary cell ranks, weights, and finite orders \(r_\alpha\),
and all residue moduli simultaneously.

\begin{proof}
Normalize in the original simplicial direction.  The normalized
complex of \(L\) is the direct sum of the coefficient modules
\((S/r_\alpha)\langle v_\alpha\rangle\) placed in their
cell dimensions \(n_\alpha\), and its differential is zero.
Exactness of normalization embeds \(NK\) in that complex,
so its differential is also zero.  On the other hand
\[
                         N\mathcal V(J)=J[1].
                                                               \tag{136.LK12}
\]
For example, normalize
\(0\to\underline J\to(J^{n+1})_n\to\mathcal V(J)\to0\):
the middle complex is the degree-one unit disk on \(J\),
and the constant subcomplex is its degree-zero copy.  This
proves (136.LK12) with its actual simplicial structure.

The map (136.LK11) is therefore just a coefficient-linear
map \((N_1K)_{\geq W_0}\to J\).  Injectivity of \(J\)
extends it to \(N_1L\).  Because \(NL\) has zero
differential, setting the other normalized components to zero
gives a chain map, hence a simplicial map
\(\lambda_j:L\to\mathcal V(J)\) extending (136.LK11).
Only actual one-dimensional cells occur in \(N_1L\).
This extension from the whole high-weight submodule handles
arbitrarily large cell weights directly; it does not leave
uncontrolled strips above each separate native weight.

For a one-dimensional cell its extension value is
\(\ell_{j\alpha}\in J^{v_\alpha}[r_\alpha]\).
Choose \(s_{j\alpha}\in S^{v_\alpha}\) such that
\(\ell_{j\alpha}=s_{j\alpha}/r_\alpha\pmod S\).
The endpoint hypothesis at that weight gives a zero-form
\(b_{j\alpha}\) on \(\Delta[1]\) with endpoints
\(0,s_{j\alpha}\).  Then
\(h_{j\alpha}=d b_{j\alpha}\) is a closed one-form
relative to the endpoints and has that integral relative period.
Replace
\[
             g_j(t_\alpha)\quad\text{by}\quad
                    g_j(t_\alpha)+h_{j\alpha}.
                                                               \tag{136.LK13}
\]
For the reference target take zero corrections.  The differential
and the full boundary of every generator are unchanged.
The actual free PD pushout therefore extends these assignments
to strict simplicial maps under \(A\).

We check the effect without reducing to a single modulus.
Represent a source two-class as \(a+\sum s_i c_i\), as
in (136.LK9).  Choose a common multiple \(N\) of its
annihilator and the finitely many relevant \(r_i\), enlarging
it if necessary to kill a remaining old two-class.  Its
one-form relation can be written
\[
                    b=b_{\mathrm{old}}-
                         \sum_i(N/r_i)s_i t_i.
                                                               \tag{136.LK14}
\]
Indeed the derivative of the required old part is
\(Na+\sum_i(N/r_i)s_i e_i\); this is closed, and its
class can be killed by enlarging \(N\) because \(H^2(A)\)
is integer torsion.  Under (136.LK13) the target image of the
two-cycle is unchanged, while (136.LK14) changes by
\(-\sum_i(N/r_i)s_i h_i\).  Formula (136.LK10) therefore
changes by exactly
\[
                        -\sum_i s_i\ell_{ji}
                              =-\lambda_j((s_i)_i).
                                                               \tag{136.LK15}
\]
The same computation is compatible with all simplicial pullbacks.
Since \(\lambda_j\) extends (136.LK11), the new linking
difference is zero throughout the specified tail.  Its equivalence
with every finite-modulus residue proves the theorem.  No choice
of a coefficient-linear family of source primitives or target
cone primitives is assumed.
\end{proof}

\subsection{The exact additional gate for reduced representables}

\paragraph{ARCH-136-LINKING-CHAIN-GATE.}
For a general coefficient-cell complex \(L\), put
\(K=\ker(L\to H^3(A))\), and suppose the natural linking
difference \(K_{\geq W}\to\mathcal V(J)\) has been
defined as above.  Write its normalized degree-one map as
\(\delta:(N_1K)_{\geq W}\to J\).  It extends to a
chain map \(NL\to J[1]\) if and only if
\[
 \delta\bigl((N_1K)_{\geq W}\cap
                      \partial(N_2L)\bigr)=0.
                                                               \tag{136.LK16}
\]
Necessity follows from the chain-map identity.  For sufficiency,
(136.LK16) defines a map from
\[
 \frac{(N_1K)_{\geq W}+\partial(N_2L)}{\partial(N_2L)}
                      \ \subseteq\ \frac{N_1L}{\partial(N_2L)}
\]
to \(J\).  Injectivity of \(J\) extends it to the right
module; compose with the quotient and set all other normalized
components to zero.  This is the required chain map.

The condition is not redundant even for an injective coefficient
target.  Take \(NL\) to be the identity disk
\(M\xrightarrow{1}M\) in degrees two and one, and
\(NK=M[1]\) its bottom subcomplex.  For
\(M=(S/r)\langle v\rangle\), where \(r\) is a nonunit
in \(R\) (for example \(r=p\) over \(\mathbb Z_{(p)}\)),
the map sending its generator
to \(T^{[v]}/r\pmod S\) is a nonzero map \(NK\to J[1]\).
It cannot extend across that identity differential.  This is
an elementary chain-level counterexample to automatic extension,
not a claimed realization by the old-algebra cell data above.

For full-boundary cell quotients, \(\partial(N_2L)=0\),
which is exactly why the preceding theorem applies.  For
\(\Delta[n]/\operatorname{sk}_0\Delta[n]\), the
normalized coefficient complex has nonzero triangle boundaries
in degree one.  A chain map on \(K\) kills
\(\partial(N_2K)\), but that subgroup can be smaller than
the intersection in (136.LK16).  Graded injectivity alone does
not remove this distinction.  Replacing the reduced-representable
attachments in ARCH-136-DEGREE-TWO-FREE by full-boundary cells
requires installing compatible boundary primitive data first;
that is a further construction problem.  Thus (136.LK16), and
not an unproved simplicial injectivity assertion, is the exact
remaining gate for that proposed application.

\paragraph{Verification status.}
The author has completed a whole saved-proof reread, including
all mixed low-degree terms, the exact cohomology sequence,
every linking-choice and all-modulus check, simplicial naturality,
the normalized coefficient extension, the corrective signs and
uniform weights, and the exact additional chain condition.
Root and reviewers15 and123 have also completed whole saved-proof
readings with PASS, including the exact chain criterion.  Their
requested nonunit clarification in its elementary illustration
has been made.  Root
contributed the chain-gate argument, and reviewer15 contributed
to the linking calculation; those are contributor readings,
not claims of independent rediscovery.  The events distinguish
the earlier communicated calculation from these saved-file reads.
